\chapter{General Project Context}

\section{Presentation of the host company}

SagemCom is an industrial group specializing in the design and manufacturing of high-value communication equipment (set-top boxes, internet boxes, smart meters). The company operates automated production lines combining electronic component assembly (SMT), functional testing, and quality control, spread across several production areas.

\subsection{Presentation of the host department}

The internship took place within the Software Engineering and Innovation Department, whose main mission is to develop internal digital tools aimed at modernizing the management and maintenance of production equipment, drawing on the industrial data generated by the manufacturing lines.

\section{Project context}

The management of industrial assets (production machines, assembly lines, test benches) traditionally relies on preventive maintenance cycles planned at fixed intervals, regardless of the actual condition of the equipment. This approach leads either to unnecessary interventions on machines still in good condition, or to unanticipated failures on machines degrading faster than expected.

The rise of industrial sensors and machine learning techniques now makes it possible to consider predictive maintenance: estimating the probability of failure, the likely failure type, the remaining useful life (RUL) and detecting abnormal machine behavior from its operating data.

\section{Problem statement}
\label{sec:problematic}

Industrial production lines rely on maintenance strategies that are either purely reactive (intervening only after a breakdown has occurred) or strictly calendar-based, where preventive tasks are performed at fixed intervals regardless of the actual condition of the equipment. Reactive maintenance leads to unplanned downtime, emergency part procurement at premium cost, and potential safety hazards for nearby operators. Calendar-based preventive maintenance, while reducing the frequency of catastrophic failures, generates unnecessary interventions on machines still in good condition and fails to catch machines degrading faster than the planned cycle.

The core problem this internship addresses is therefore threefold:
\begin{enumerate}
\item \textbf{Absence of condition-based visibility.} Without continuous analysis of sensor data (temperature, rotational speed, torque, tool wear), maintenance teams have no objective basis for deciding \emph{when} a specific machine actually needs attention, leading to either over-maintenance or surprise failures.
\item \textbf{Fragmented tooling.} Enterprise Asset Management (EAM) systems, when they exist, typically handle only the administrative workflow (work orders, intervention scheduling, spare-parts inventory) and offer no integration with predictive analytics or field-level decision support. Machine Learning models, when developed separately, remain disconnected from the operational cycle and are never used by the technicians who need them.
\item \textbf{Lack of contextualized field support.} Technicians in the field must currently consult separate paper manuals, internal wikis, or call the head office to interpret an alert or decide on a corrective action. There is no single interface that combines the machine's real-time health status, its maintenance history, and the relevant documentation into a single, actionable answer.
\end{enumerate}

Bridging these three gaps (condition-based visibility, integrated predictive analytics, and contextualized field support) within a single, cohesive platform is the problem this project set out to solve.

\section{Internship objectives}
\label{sec:internship-objectives}

The main objective of this internship is to design and implement a single, integrated platform that combines the following capabilities:

\begin{itemize}
\item \textbf{EAM operational management:} management of users, machines, work orders, intervention orders, and maintenance scheduling.
\item \textbf{Predictive maintenance:} design and integration of machine learning models (P1 to P7) covering failure probability prediction, failure type classification, remaining useful life estimation, anomaly detection, intervention prioritization and scheduling, and spare parts coordination.
\item \textbf{Intelligent support for field teams:} implementation of contextualized alerts, prediction explainability to answer ``why this alert?'', and a conversational assistant capable of answering technicians' questions using maintenance documentation and the real-time health status of machines.
\end{itemize}

To achieve these objectives, the platform development will focus on:

\begin{itemize}
\item \textbf{Backend and frontend development} of the EAM platform.
\item \textbf{Design and training of the P1 to P7 predictive maintenance models.}
\item \textbf{Integration of a Retrieval-Augmented Generation (RAG) service} for documentary assistance, connected to real-time machine learning predictions.
\item \textbf{Validation of the complete solution} through functional testing and performance evaluation of the developed models.
\end{itemize}

\section{Working methodology}
\label{sec:methodology}

The internship followed an iterative, sprint-based methodology inspired by Scrum, adapted to the constraints of a solo-developer end-of-studies project within an industrial environment. Each feature was the subject of a written specification (requirements document), followed by a detailed implementation plan, before being built and verified. This methodology made it possible to maintain full traceability of design decisions throughout the six months of the internship.

\subsection{Sprint-based development}

Work was organised into two-week sprints, each ending with a demonstrable milestone. A typical sprint followed four phases:

\begin{enumerate}
\item \textbf{Sprint planning:} the objectives for the coming two weeks were defined, broken down into concrete tasks with estimated effort, and prioritised against the project roadmap.
\item \textbf{Development and integration:} features were implemented, tested locally, and merged into the working branch. Dependencies between frontend, backend, and ML components were resolved incrementally.
\item \textbf{Sprint review:} the results were demonstrated to the internship supervisor and, when relevant, to the company's technical stakeholders (Head Technician, Administrator). Feedback was recorded and fed into the next planning cycle.
\item \textbf{Sprint retrospective:} what worked well and what needed adjustment in the process was identified, leading to concrete improvements for the next iteration (tooling, documentation practices, testing strategy).
\end{enumerate}

\subsection{Internship lifecycle}

The six-month internship was structured into five sequential phases, each building on the previous one. Figure~\ref{fig:planning} provides a visual timeline.

\paragraph{Phase 1: Analysis and architecture (month 1).}
The existing EAM system at SagemCom was studied in depth: business processes, data model, user roles, and pain points. The four-component architecture (frontend, backend, ML microservice, RAG service) was designed during this phase, along with the initial data model and API contracts.

\paragraph{Phase 2: Core platform development (months 1--2).}
The backend services (user management, machine registry, work orders, intervention orders) and the frontend dashboards for the three roles (Administrator, Head Technician, Technician) were built incrementally. Each feature followed the same loop: specification, implementation plan, build, verification.

\paragraph{Phase 3: Machine Learning integration (months 2--4).}
The ML microservice was developed progressively, starting with the six initially scoped prediction problems (P1--P6) and expanding to the Wave~2 signal-fusion layer and the spare-parts demand forecasting module (P7). Each model followed an eight-step pipeline (define, collect, prepare, split, choose, train, evaluate, tune) applied identically across all problems. This phase also included the RAG service for document ingestion and the ML--RAG bridge for conversational assistance.

\paragraph{Phase 4: Explainability, alerts, and production readiness (months 4--5).}
Explainability panels, role-aware alert routing, and the readiness score were added. The system was deployed on a Kubernetes cluster (Azure AKS) with Terraform-managed infrastructure and Helm-based release management.

\paragraph{Phase 5: Testing and validation (month 6).}
Comprehensive testing was carried out: unit tests for business logic, integration tests for API endpoints, and end-to-end validation of the ML pipeline against known failure scenarios. The label-leakage and validation-methodology issues identified during this phase led to significant revisions in the model training approach (Chapter~5).

\subsection{Traceability and documentation}

Every design decision was recorded in a structured format: the problem being solved, the alternatives considered, the choice made, and the justification. This documentation served two purposes: it made it possible to revisit early assumptions when later evidence contradicted them (as happened with the label-leakage issue), and it provided the foundation for this report. The internship schedule shown in Figure~\ref{fig:planning} compresses this six-month process into a single Gantt chart, with overlapping bars reflecting the parallel nature of the work.

\section{Internship schedule}

\definecolor{gantt1}{HTML}{274472}
\definecolor{gantt2}{HTML}{34507B}
\definecolor{gantt3}{HTML}{415C85}
\definecolor{gantt4}{HTML}{4E688E}
\definecolor{gantt5}{HTML}{5B7498}
\definecolor{gantt6}{HTML}{6880A1}
\definecolor{gantt7}{HTML}{758CAA}
\definecolor{gantt8}{HTML}{8298B4}
\definecolor{gantt9}{HTML}{8EA4BD}
\definecolor{gantt10}{HTML}{9BB0C7}
\definecolor{gantt11}{HTML}{A8BCD0}
\definecolor{gantt12}{HTML}{B5C8D9}
\definecolor{gantt13}{HTML}{C2D4E3}
\definecolor{gantt14}{HTML}{CFE0EC}

\begin{figure}[p]
\noindent\makebox[\textwidth][c]{%
\begin{ganttchart}[
    vgrid,
    hgrid,
    x unit=1.35cm,
    y unit title=1.1cm,
    y unit chart=1.05cm,
    title height=1,
    bar height=0.6,
    bar label font=\small,
    bar label node/.append style={text width=4.6cm, align=right},
    title label font=\bfseries\normalsize,
    title/.style={draw=none, fill=black!10},
]{1}{7}
\gantttitle{Internship schedule (months)}{7} \\
\gantttitlelist{1,...,7}{1} \\
\ganttbar[bar/.append style={fill=gantt1}]{Analysis \& architecture}{1.0}{1.4} \\
\ganttbar[bar/.append style={fill=gantt2}]{Backend (users/machines)}{1.1}{1.9} \\
\ganttbar[bar/.append style={fill=gantt3}]{Backend (work/interv. orders)}{1.6}{2.3} \\
\ganttbar[bar/.append style={fill=gantt4}]{Frontend (dashboards)}{1.8}{2.5} \\
\ganttbar[bar/.append style={fill=gantt5}]{ML P1 to P3}{2.2}{3.2} \\
\ganttbar[bar/.append style={fill=gantt6}]{ML P4 to P6}{2.8}{3.7} \\
\ganttbar[bar/.append style={fill=gantt7}]{DST fusion (Wave 2)}{3.3}{3.9} \\
\ganttbar[bar/.append style={fill=gantt8}]{RAG (ingestion/S3)}{3.5}{4.2} \\
\ganttbar[bar/.append style={fill=gantt9}]{ML-RAG bridge}{3.9}{4.6} \\
\ganttbar[bar/.append style={fill=gantt10}]{P7 (spare parts)}{4.2}{5.1} \\
\ganttbar[bar/.append style={fill=gantt11}]{Explainability \& alerts}{4.8}{5.6} \\
\ganttbar[bar/.append style={fill=gantt12}]{Phase 4 (model redesign)}{5.4}{6.4} \\
\ganttbar[bar/.append style={fill=gantt13}]{Testing \& fixes}{6.0}{6.8} \\
\ganttbar[bar/.append style={fill=gantt14}]{Report writing}{6.5}{7.0}
\end{ganttchart}%
}
\caption{End-of-studies internship schedule}
\label{fig:planning}
\end{figure}
